\chapter{Referensi dan Sumber Belajar}
\label{app:referensi-tambahan}

Seluruh bahan kajian, contoh, praktikum, dan latihan pada buku ini
diturunkan dari dua buku sumber lokal. Kedua buku tersebut menjadi
acuan tunggal agar isi buku tetap sejalan dengan RPS dan tidak
memasukkan topik di luar cakupan mata kuliah.

\section{Buku Sumber Utama}

\begin{enumerate}
    \item \textbf{Indra Gunawan. (2019).}
          \textit{Pemrograman Bahasa Assembly Mikroprosesor 8086}.
          Guepedia. \cite{gunawan2019}

          Buku ini membahas perkembangan mikroprosesor, komponen dasar
          8086, interupsi, persiapan lingkungan assembly, transfer data,
          pencetakan karakter dan kalimat, operasi logika, tipe data,
          mode pengalamatan, array, percabangan dan perulangan, stack,
          prosedur dan macro, operasi aritmatika, operasi file, string,
          grafik, serta layanan interupsi yang sering digunakan.

    \item \textbf{S'to. (2001).}
          \textit{Pemrograman dengan Bahasa Assembly, Edisi Online
          Versi 1.0}. IlmuKomputer.Com. \cite{sto2001}

          Buku ini membahas sistem bilangan, organisasi memori,
          interupsi, register, memulai assembly, program COM, pencetakan
          huruf dan kalimat, operasi aritmatika, pointer, manipulasi bit
          dan logika, mode pengalamatan, perbandingan dan lompatan,
          stack, masukan keyboard, prosedur, macro, operasi layar,
          operasi string, pencetakan angka, pengaksesan port, program
          berparameter, operasi file, program residen, grafik, program
          EXE, dan Turbo Debugger.
\end{enumerate}

\section{Peta Bahan Kajian ke Sumber}

Tabel~\ref{tab:peta-sumber} menunjukkan bab buku ini, topik utamanya,
serta bagian dan halaman yang menjadi rujukan pada kedua buku sumber.
Nomor halaman mengikuti penanda halaman berkas \texttt{.tex} sumber.

\begin{table}[htbp]
\centering
\caption{Peta bahan kajian ke dua buku sumber}
\label{tab:peta-sumber}
\small
\begin{tabular}{@{}p{1.5cm}p{4.3cm}p{4.2cm}p{4.3cm}@{}}
\toprule
\textbf{Bab} & \textbf{Topik} & \textbf{Sumber [1]} &
\textbf{Sumber [2]} \\
\midrule
III & Sistem bilangan & Bagian 1, hlm. 10--11 &
BAB I, hlm. 6--8 \\
\addlinespace
IV & BIU dan EU & Bagian 2, hlm. 12--18 &
BAB II, hlm. 9--10 \\
\addlinespace
V & Register dan flag & Bagian 2, hlm. 13--18 &
BAB IV, hlm. 17--20 \\
\addlinespace
VI & Segmentasi memori & Bagian 2.1.1, hlm. 13--15 &
BAB II, hlm. 9--13 \\
\addlinespace
VII & Memulai assembly & Bagian 4, hlm. 26--34 &
BAB V--VI, hlm. 21--26 \\
\addlinespace
VIII & Transfer data & Bagian 5, hlm. 35--44 &
BAB V, hlm. 21--24 \\
\addlinespace
IX & Interupsi BIOS/DOS & Bagian 3, hlm. 19--25; Bagian 18,
hlm. 179--204 & BAB III, hlm. 14--16; BAB VII, hlm. 27--31; BAB XII,
hlm. 63--67; BAB XV, hlm. 78--84 \\
\addlinespace
X & Aritmatika & Bagian 14, hlm. 127--142 &
BAB VIII, hlm. 32--40 \\
\addlinespace
XI & Logika dan shift & Bagian 7, hlm. 68--73 &
BAB X, hlm. 47--51 \\
\addlinespace
XII & Mode pengalamatan & Bagian 9, hlm. 78--92 &
BAB XI, hlm. 52--62 \\
\addlinespace
XIII & Percabangan dan perulangan & Bagian 11, hlm. 101--114 &
BAB XIII, hlm. 68--72 \\
\addlinespace
XIV & Stack, prosedur, macro & Bagian 12--13, hlm. 115--126 &
BAB XIV, hlm. 73--77; BAB XVI--XVII, hlm. 85--92 \\
\addlinespace
XV & Array, string, bilangan & Bagian 10, hlm. 93--100; Bagian 14.1,
hlm. 127--131; Bagian 16, hlm. 154--158 & BAB IX, hlm. 41--46;
BAB XIX--XX, hlm. 110--121 \\
\addlinespace
XVI & Debugging & Bagian 4.2--4.3, hlm. 29--32 &
BAB XXVII, hlm. 194--202 \\
\bottomrule
\end{tabular}
\end{table}

\section{Perangkat Lunak Pendukung}

Praktik pada buku ini memakai perangkat yang dibahas pada kedua buku
sumber, yaitu:

\begin{itemize}
    \item \textbf{TASM} untuk meng-assemble berkas \texttt{.ASM} menjadi
          \texttt{.OBJ};
    \item \textbf{TLINK} untuk men-link berkas \texttt{.OBJ} menjadi
          \texttt{.EXE} atau \texttt{.COM};
    \item \textbf{DEBUG} untuk memeriksa register, memori, dan
          menjalankan program per instruksi;
    \item \textbf{emu8086} untuk menulis, menjalankan, dan men-debug
          program;
    \item \textbf{DOSBox} untuk menyediakan lingkungan DOS bagi TASM,
          TLINK, dan DEBUG;
    \item \textbf{Turbo Debugger} untuk menelusuri program per instruksi
          dan memeriksa register, flag, memori, serta stack.
\end{itemize}

\section{Cara Menggunakan Lampiran Ini}

\begin{enumerate}
    \item Gunakan buku ini sebagai kerangka utama dan penentu urutan
          belajar.
    \item Ketika menemukan istilah yang belum jelas, buka
          Lampiran~\ref{app:glosarium}.
    \item Ketika memerlukan penjelasan tambahan suatu topik, buka bagian
          dan halaman yang tercantum pada Tabel~\ref{tab:peta-sumber}.
    \item Bandingkan penjelasan pada kedua sumber: buku [1] menekankan
          contoh program dan langkah praktik, sedangkan buku [2]
          menekankan penjelasan bertahap dan latihan.
    \item Setelah menyelesaikan seluruh bab, gunakan
          Lampiran~\ref{app:rubrik} untuk memeriksa ketercapaian
          kompetensi.
\end{enumerate}

\begin{infobox}
\textbf{Catatan rujukan.} Buku ini hanya memakai dua buku sumber di
atas. Sitasi \cite{gunawan2019} dan \cite{sto2001} merujuk pada kedua
buku tersebut. Tidak ada tautan atau sumber daring yang digunakan.
\end{infobox}
